\documentclass[10pt, oneside]{article}
\usepackage{amsmath, amsthm, amssymb, calrsfs, wasysym, verbatim, bbm, color, graphics, geometry}
\usepackage{graphicx}
\graphicspath{ {./assets/} }
\geometry{tmargin=.75in, bmargin=.75in, lmargin=.75in, rmargin = .75in}

\newcommand{\R}{\mathbb{R}}
\newcommand{\C}{\mathbb{C}}
\newcommand{\Z}{\mathbb{Z}}
\newcommand{\N}{\mathbb{N}}
\newcommand{\Q}{\mathbb{Q}}
\newcommand{\Cdot}{\boldsymbol{\cdot}}

\newtheorem{thm}{Theorem}
\newtheorem{defn}{Definition}
\newtheorem{conv}{Convention}
\newtheorem{rem}{Remark}
\newtheorem{lem}{Lemma}
\newtheorem{cor}{Corollary}
\usepackage[backend=biber,style=authoryear]{biblatex}
\addbibresource{reference.bib}

\title{Summary Report: Candidate literary paper to be used}
\author{Alejandro Ouslan}
\date{2026-02-08}

\begin{document}

\maketitle
\tableofcontents

\vspace{.25in}

\section{Paper 1: On the four types of weight functions for spatial contiguity matrix}

\subsection{Summary}
The studies of spatial complexity have resulted in new understanding of the distance-decay functions of geography.
The weight functions of spatial autocorrelation should be consistent with the functions of spatial distribution and
interaction. This paper is on “spatial autocorrelation analysis” of spatial autocorrelation. The 1-dimension spatial
autocorrelation function (ACF) and partial autocorrelation function (PACF) of space series are employed to analyze four
kinds of weight function in common use for the 2-dimensional spatial autocorrelation model. The aim of this study is
at how to select a proper weight function to construct a spatial contiguity matrix for spatial analysis. The scopes
of application of different weight functions are defined in terms of the characters of their ACFs and PACFs. \cite{chen2012four}

\subsection{Metadata}
\begin{enumerate}
	\item \textbf{Link:} https://link.springer.com/article/10.1007/s12076-011-0076-6
	\item \textbf{Journal Name:} Letters in Spatial and Resource Sciences
	\item \textbf{Rank:} Q2
\end{enumerate}

\subsection{Conclusions}
\begin{itemize}
	\item \textbf{Objective:} To provide an objective criterion for selecting spatial weight functions based on the consistency between spatial distribution, interaction, and scale.
	\item \textbf{Methodology:} The author uses 1-dimensional Autocorrelation Functions (ACF) and Partial Autocorrelation Functions (PACF) to analyze the behavior and "spatial memory" of four common weight functions used in 2-dimensional spatial analysis.
	\item \textbf{Key Findings:}
	      \begin{itemize}
		      \item Weight functions must match the underlying geographical process (e.g., power-law vs. exponential distributions).
		      \item The choice of function is critical at "critical states" (the transition between chaos and order), where different functions yield significantly different analytical results.
	      \end{itemize}
	\item \textbf{Conclusion:}
	      \begin{itemize}
		      \item \textbf{Power functions} are best for large-scale complex systems.
		      \item \textbf{Exponential functions} suit medium-to-small scale simple systems.
		      \item \textbf{Step functions} are ideal for local-level analysis and negative correlation patterns.
	      \end{itemize}
\end{itemize}

\section{Template}

\subsection{Summary}
\cite{chen2012four}

\subsection{Metadata}
\begin{enumerate}
	\item \textbf{Link:}
	\item \textbf{Journal Name:}
	\item \textbf{Rank:}
\end{enumerate}

\subsection{Conclusions}


\printbibliography
\end{document}
